\DocumentMetadata{
  pdfversion=2.0,
  pdfstandard=ua-2,
  lang=en-US,
  tagging=on,
  tagging-setup={math/setup=mathml-SE,tabsorder=structure}
}
\documentclass[11pt,letterpaper]{article}
\usepackage[margin=0.68in,includefoot]{geometry}
\usepackage{newcomputermodern}
\usepackage{amsmath,amscd,mathtools}
\providecommand{\square}{\mdlgwhtsquare}
\usepackage[hidelinks]{hyperref}
\hypersetup{
  pdftitle={Algebra First-Year Exam, May 2024, Part 1},
  pdfauthor={University of Florida Department of Mathematics},
  pdfsubject={Graduate mathematics examination},
  pdfdisplaydoctitle=true,
  bookmarksopen=true
}
\setcounter{secnumdepth}{0}
\setlength{\parindent}{0pt}
\setlength{\parskip}{0.28em}
\setlength{\emergencystretch}{3em}
\setlength{\leftmargini}{2.15em}
\setlength{\leftmarginii}{2.65em}
\setlength{\leftmarginiii}{3em}
\renewcommand{\labelenumi}{\textbf{\theenumi.}}
\pagestyle{empty}
\raggedbottom
\allowdisplaybreaks
\newenvironment{examproblems}[1]{%
  \begin{enumerate}%
  \setcounter{enumi}{#1}%
  \setlength{\topsep}{0.35em}%
  \setlength{\partopsep}{0pt}%
  \setlength{\itemsep}{0.48em}%
  \setlength{\parsep}{0.18em}%
}{\end{enumerate}}
\newenvironment{examsubparts}{%
  \begin{enumerate}%
  \setlength{\topsep}{0.18em}%
  \setlength{\partopsep}{0pt}%
  \setlength{\itemsep}{0.16em}%
  \setlength{\parsep}{0.1em}%
}{\end{enumerate}}
\newenvironment{examinstructions}{%
  \par\begingroup\itshape
}{%
  \par\endgroup\vspace{0.18em}
}
\makeatletter
\renewcommand\section{\@startsection{section}{1}{0pt}%
  {-1.25ex plus -.25ex minus -.1ex}%
  {0.55ex plus .1ex}%
  {\normalfont\large\bfseries}}
\renewcommand{\@maketitle}{%
  \newpage\null\vskip -1.2em%
  {\footnotesize\bfseries
    \href{https://www.ufl.edu/}{University of Florida}, \href{https://math.ufl.edu/}{Department of Mathematics}\par}%
  \vskip 0.18em%
  \tagstructbegin{tag=Title}%
    {\LARGE\bfseries\href{https://gma.math.ufl.edu/exams/algebra-first-year/}{Algebra First-Year Exam}, May 2024, Part 1\par}%
  \tagstructend%
  \vskip 0.35em%
  \tagmcbegin{artifact}%
    \hrule height 1pt%
  \tagmcend%
  \vskip 0.55em%
}
\makeatother
\title{Algebra First-Year Exam, May 2024, Part 1}
\author{}
\date{}
\begin{document}
\maketitle
\thispagestyle{empty}
\begin{examinstructions}
Answer four problems. (If you turn in more, the first four will be graded.) Within reason, you may quote theorems as long as you state them clearly.
\end{examinstructions}
\begin{examproblems}{0}
\item
State and prove Lagrange’s Theorem.
\item
Let~\(G\) be a finite group, and let~\(H\) and~\(J\) be normal subgroups of~\(G\).
\begin{examsubparts}
\item[\textnormal{(a)}]
Prove that \(H \cap J\) is a normal subgroup of~\(G\);
\item[\textnormal{(b)}]
Prove that if \(H \cap J = 1\) then, for every \(h \in H\) and every \(j \in J\), we have \(hj = jh\).
\end{examsubparts}
\item
Let~\(G\) be a group acting on a non empty set~\(S\), and~\(H\) be a normal subgroup of~\(G\). Assume that for any \(x_1, x_2 \in S\) there is a unique \(h \in H\) such that \(h(x_1) = x_2\). For \(x \in S\) let \(G_x\) be the stabilizer of~\(x\). Show that for any \(x \in S\), \(G\) is a semi-direct product of \(G_x\) and~\(H\). Do not assume that~\(G\) is finite.
\item
This problem has two parts.
\begin{examsubparts}
\item[\textnormal{(a)}]
Prove that a group~\(G\) of order \(245\) has a unique subgroup~\(H\) of order \(49\).
\item[\textnormal{(b)}]
Show that if the above subgroup~\(H\) is cyclic then~\(G\) is abelian.
\end{examsubparts}
\item
Give the definition of a solvable (finite) group. (You may give any of the equivalent definitions.) Prove, using your definition, that the symmetric group on four letters is solvable.
\item
This problem has two parts.
\begin{examsubparts}
\item[\textnormal{(a)}]
Find all the conjugacy classes of elements of order~\(2\) in \(S_7\).
\item[\textnormal{(b)}]
Do the same for \(A_7\). In both cases, justify your answer.
\end{examsubparts}
\end{examproblems}
\end{document}
